% Part: sets-functions-relations
% Chapter: sets
% Section: basics

\documentclass[../../../include/open-logic-section]{subfiles}

\begin{document}

\olfileid[fa-IR]{sfr}{set}{bas}
\olsection{اصل گسترش}

یک \emph{مجموعه} گردایه‌ای از اشیاست که {به‌منزلهٔ} شیئی واحد در
نظر گرفته می‌شود. اشیای سازندهٔ مجموعه را \emph{عناصر} یا
\emph{اعضای} آن می‌نامند. اگر $x$ !!a{element} مجموعهٔ~$A$ باشد،
می‌نویسیم $x \in A$؛ و اگر چنین نباشد، می‌نویسیم $x \notin A$.
مجموعه‌ای را که فاقد !!{element}s است، \emph{مجموعهٔ تهی} می‌نامند و
آن را با~``$\emptyset$'' نشان می‌دهند.

\begin{explain}
شیوهٔ \emph{مشخص‌کردن} مجموعه، \emph{ترتیب} !!{element}s آن، و حتی
اینکه در شمارش !!{element}s آن هر کدام را \emph{چند بار} بشماریم،
هیچ‌یک اهمیتی ندارد. آنچه اهمیت دارد این است که !!{element}s آن چه
چیزهایی‌اند. این نکته را در اصل زیر صورت‌بندی می‌کنیم.
\end{explain}

\begin{defn}[اصل گسترش]
  اگر $A$ و $B$ مجموعه باشند، آنگاه $A = B$ اگر و تنها اگر
  هر !!{element} مجموعهٔ~$A$، !!a{element} مجموعهٔ~$B$ نیز باشد، و
  برعکس.
\end{defn}

اصل گسترش به ما اجازه می‌دهد از برخی نمادگذاری‌ها استفاده کنیم.
به‌طور کلی، اگر اشیایی چون $a_{1}$، \dots، $a_{n}$ داشته باشیم، آنگاه
$\{a_{1}, \dots, a_{n}\}$ \emph{یگانه} مجموعه‌ای است که !!{element}s آن
$a_1, \ldots, a_n$ هستند. بر واژهٔ ``\emph{یگانه}'' تأکید می‌کنیم،
زیرا اصل گسترش می‌گوید تنها \emph{یک} چنین مجموعه‌ای می‌تواند وجود
داشته باشد. در واقع، اصل گسترش برابری‌های زیر را نیز مجاز می‌سازد:
  \[
    \{a, a, b\} = \{a, b\} = \{b,a\}.
  \]
این برابری‌ها همان نکته را بیان می‌کنند که هنگام بررسی مجموعه‌ها،
ترتیب !!{element}s آنها یا تعداد دفعاتی که مشخص شده‌اند، اهمیتی ندارد.

\begin{tagblock}{novice}
\begin{ex}
هرگاه تعدادی شیء داشته باشید، می‌توانید آنها را در یک مجموعه گرد آورید.
برای نمونه، مجموعهٔ خواهران و برادران ریچارد مجموعه‌ای است که یک نفر را
در بر می‌گیرد و می‌توان آن را به‌صورت $S=\{\textrm{روث}\}$ نوشت.
مجموعهٔ اعداد صحیح مثبت کوچک‌تر از $4$ برابر $\{1, 2, 3\}$ است، اما
می‌توان آن را به‌صورت $\{3, 2, 1\}$ یا حتی $\{1, 2, 1, 2, 3\}$ نیز نوشت.
اینها بنا بر اصل گسترش همگی یک مجموعه‌اند. زیرا هر !!{element}
$\{1, 2, 3\}$، !!a{element} $\{3, 2, 1\}$ (و $\{1,
2, 1, 2, 3\}$) نیز هست، و برعکس.
\end{ex}
\end{tagblock}

اغلب مجموعه‌ای را به‌وسیلهٔ ویژگی مشترک !!{element}s آن مشخص می‌کنیم.
برای این کار از نمادگذاری اختصاری زیر استفاده می‌کنیم:
$\Setabs{x}{\phi(x)}$، که در آن $\phi(x)$ نمایانگر ویژگی‌ای است
که~$x$ باید داشته باشد تا در شمار !!{element}s مجموعه قرار گیرد.

\begin{tagblock}{novice}
\begin{ex}
در مثال ما می‌شد $S$ را به‌صورت زیر نیز مشخص کرد:
\[
S = \Setabs{x}{x \text{ خواهر یا برادر ریچارد است}}.
\]
\end{ex}
\end{tagblock}

\begin{tagblock}{math}
\begin{ex}
عددی را \emph{کامل} می‌نامیم اگر و تنها اگر با مجموع مقسوم‌علیه‌های
سرهٔ خود برابر باشد (یعنی عددهایی که آن را بدون باقیمانده تقسیم می‌کنند
اما با خود آن عدد برابر نیستند). برای نمونه، $6$ کامل است زیرا
مقسوم‌علیه‌های سرهٔ آن $1$، $2$ و~$3$ هستند و $6 = 1 + 2 + 3$.
در واقع، $6$ تنها عدد صحیح مثبت کوچک‌تر از $10$ است که کامل است.
ازاین‌رو، با استفاده از اصل گسترش می‌توانیم بگوییم:
\[
	\{6\} = \Setabs{x}{x\text{ کامل است و }0 \leq x \leq 10}
\]
نماد سمت راست را چنین می‌خوانیم: ``مجموعهٔ $x$هایی که $x$
کامل است و $0 \leq x \leq 10$''. این همانی تأیید می‌کند که هنگام
بررسی مجموعه‌ها، شیوهٔ مشخص‌کردن آنها اهمیتی ندارد. به‌طور کلی‌تر،
اصل گسترش تضمین می‌کند که همواره تنها یک مجموعه از $x$هایی وجود دارد
که $\phi(x)$ درباره‌شان صادق است. پس اصل گسترش مجوز می‌دهد
$\Setabs{x}{\phi(x)}$ را \emph{یگانه} مجموعهٔ $x$هایی بنامیم
که~$\phi(x)$ درباره‌شان صادق است.
\end{ex}
\end{tagblock}

اصل گسترش راهی برای اثبات یکسان‌بودن مجموعه‌ها به دست می‌دهد: برای
اثبات $A = B$، نشان دهید هرگاه $x \in A$ باشد، $x \in B$ نیز هست،
و هرگاه $y \in B$ باشد، $y \in A$ نیز هست.

\begin{prob}
ثابت کنید حداکثر یک مجموعهٔ تهی وجود دارد؛ یعنی نشان دهید اگر $A$ و $B$
مجموعه‌هایی فاقد !!{element}s باشند، آنگاه $A = B$.
\end{prob}

\end{document}
